% !TeX program = lualatex
% ================================================================
% Polish Tier 3 Vocabulary Version of SequencesNthTermIdeasStarters
%
% Generated by the server using the following settings:
%
%   language            Polish (pol)
%   text direction      left to right
%   font                Noto Serif
%   fallback font       Noto Serif
%   built from          latex/starters/targeted/KS3/sequences/sequencesNthTermIdeasStarters.tex
%   vocabulary key      latex/starters/targeted/KS3/sequences/sequencesNthTermIdeasStarters/L2/pol/sequencesNthTermIdeasStarters_polishKey.tex
%   tier 3 vocabulary   green   RGB 0 128 0
%   English text        black   RGB 0 0 0
%   Polish text         purple  RGB 112 48 160
%   layout macros       version 3
%
% Please don't edit any information in this comment block - the server
% compares it against the current settings and rebuilds this file when
% they differ, so changes here would be overwritten. However, feel free
% to make updates to the LaTeX below if you want to edit it.
% ================================================================

\documentclass{beamer}
% ================================================================
% Copied from the English sheet this was translated from, so that
% anything it sets up for itself still works here
% ================================================================
\usepackage{tikz}
\usepackage{amsmath}

% ---- Minimal styling ----
\setbeamertemplate{navigation symbols}{}
\setbeamertemplate{footline}{}
\setbeamercolor{background canvas}{bg=white}
\setbeamercolor{frametitle}{fg=black}
\setbeamerfont{frametitle}{size=\Large, series=\bfseries}
\usefonttheme{professionalfonts}

% ---- Global tikz baseline ----
\tikzset{every picture/.style={baseline=(current bounding box.south)}}

% ================================================================
% TikZ pattern commands
% ================================================================

% Matchstick squares: n squares in a row
\newcommand{\matchsquares}[2][black]{%
  \begin{tikzpicture}[scale=0.60]
    \pgfmathsetmacro\n{#2}
    \color{#1}
    \foreach \x in {0,...,\n} {
      \draw[very thick, cap=round] (\x,0) -- (\x,1);
    }
    \foreach \x in {0,...,\the\numexpr\n-1\relax} {
      \draw[very thick, cap=round] (\x,0)   -- (\x+1,0);
      \draw[very thick, cap=round] (\x,1)   -- (\x+1,1);
    }
  \end{tikzpicture}%
}

% Matchstick triangles: chain of n triangles
\newcommand{\matchtriangles}[2][black]{%
  \begin{tikzpicture}[scale=0.60]
    \pgfmathsetmacro\n{#2}
    \pgfmathsetmacro\h{0.866}
    \color{#1}
    \draw[very thick] (0,0) -- (\n,0);
    \draw[very thick] (0,0)
      \foreach \k in {1,...,\n} {
        -- ({\k - 0.5}, {\h}) -- (\k,0)
      };
  \end{tikzpicture}%
}

% Rectangular dots: n columns x (n+1) rows
\newcommand{\rectdots}[2][black]{%
  \begin{tikzpicture}[scale=0.60]
    \pgfmathsetmacro\n{#2}
    \color{#1}
    \foreach \i in {0,...,\the\numexpr\n-1\relax} {
      \foreach \j in {0,...,\n} {
        \fill (0.4*\i, 0.4*\j) circle (0.07);
      }
    }
  \end{tikzpicture}%
}

% Cross / plus dots
\newcommand{\crossdots}[2][black]{%
  \begin{tikzpicture}[scale=0.60]
    \pgfmathsetmacro\n{#2}
    \color{#1}
    \fill (0,0) circle (0.07);
    \ifnum\n>1
      \foreach \d in {1,...,\the\numexpr\n-1\relax} {
        \fill (0,   0.4*\d) circle (0.07);
        \fill (0,  -0.4*\d) circle (0.07);
        \fill ( 0.4*\d, 0)  circle (0.07);
        \fill (-0.4*\d, 0)  circle (0.07);
      }
    \fi
  \end{tikzpicture}%
}

% 4-times-table dots: n rows of 4 dots
\newcommand{\fourtabledots}[2][black]{%
  \begin{tikzpicture}[scale=0.60]
    \pgfmathsetmacro\n{#2}
    \color{#1}
    \foreach \row in {0,...,\the\numexpr\n-1\relax} {
      \foreach \col in {0,1,2,3} {
        \fill (0.4*\col, 0.4*\row) circle (0.07);
      }
    }
  \end{tikzpicture}%
}

% ================================================================
% Answer-overlay helpers
% ================================================================
\newcommand{\ablank}[1]{%
  \alt<2>{\textcolor{red}{#1}}{\underline{\phantom{#1}}}%
}
\newcommand{\ashow}[1]{\uncover<2->{\textcolor{red}{\small #1}}}
\newcommand{\ashowq}[1]{\alt<2->{\textcolor{red}{\small #1}}{?}}

% Answers drawn straight on to a TikZ diagram, revealed on overlay 2 like \ashow.
% Space is reserved on overlay 1, so the diagram does not move between slides.
\usetikzlibrary{overlay-beamer-styles}
\tikzset{
  ans/.style     = {red, thick, visible on=<2->},
  ansfill/.style = {red!15, visible on=<2->},
  anslab/.style  = {red, font=\tiny, visible on=<2->},
}

% ================================================================

% Characters this language's font has not got are borrowed from another,
% rather than being silently dropped - see docs/translated-sheet-latex.md
\directlua{%
  if luaotfload and luaotfload.add_fallback then
    luaotfload.add_fallback("eallatinfallback",
      { "Noto Serif:mode=harf;" })
  else
    tex.error("this luaotfload is too old to register a fallback font")
  end
}
\usepackage[english]{babel}
\babelprovide[import]{polish}
\babelfont[polish]{rm}[Renderer=Harfbuzz, BoldFont={Noto Serif}, ItalicFont={Noto Serif}, BoldItalicFont={Noto Serif}, RawFeature={fallback=eallatinfallback}]{Noto Serif}
\babelfont[polish]{sf}[Renderer=Harfbuzz, BoldFont={Noto Serif}, ItalicFont={Noto Serif}, BoldItalicFont={Noto Serif}, RawFeature={fallback=eallatinfallback}]{Noto Serif}
\newcommand{\ealtext}[1]{\foreignlanguage{polish}{#1}}
\newcommand{\ealtextblock}[1]{{\raggedright\foreignlanguage{polish}{#1}\par}}
% ================================================================
% EAL layout helpers
% ================================================================
\usepackage{xcolor}

\definecolor{ealkeycolour}{RGB}{0,128,0}
\definecolor{ealenglish}{RGB}{0,0,0}
\definecolor{ealtwocolour}{RGB}{112,48,160}

\newsavebox{\ealboxen}
\newsavebox{\ealboxtr}
\newlength{\ealglosswd}
\newlength{\ealglossraise}
\setlength{\ealglossraise}{1.15em}

% a tier 3 word, in English and in translation alike
\newcommand{\ealkey}[1]{\textcolor{ealkeycolour}{#1}}
\newcommand{\ealkeytr}[1]{\textcolor{ealkeycolour}{\ealtext{#1}}}

% One translated word above one English word. Built from kernel
% primitives, never a tabular, which would break wherever the sheet
% loads array, booktabs or tabularx. The box is as wide as the wider of
% the two, so a long translation pushes its neighbours aside instead of
% overlapping them, and it claims its full height so a table row or a
% TikZ node makes room for it.
\newcommand{\ealgloss}[2]{%
  \sbox{\ealboxen}{\ealkey{#1}}%
  \sbox{\ealboxtr}{\small\textcolor{ealtwocolour}{\ealtext{#2}}}%
  \setlength{\ealglosswd}{\wd\ealboxen}%
  \ifdim\wd\ealboxtr>\ealglosswd\setlength{\ealglosswd}{\wd\ealboxtr}\fi
  \leavevmode
  \makebox[\ealglosswd][c]{%
    \makebox[0pt][c]{\usebox{\ealboxen}}%
    \makebox[0pt][c]{\raisebox{\ealglossraise}{\usebox{\ealboxtr}}}%
  }%
}

% opens up line spacing around a block of glosses, so the raised
% translations sit clear of the line above
\newenvironment{ealglossed}{\par\linespread{2.1}\selectfont\ignorespaces}{\par}

% a whole translated sentence above its English counterpart. block level,
% so it does not belong inside a tabular cell
\newcommand{\ealpara}[2]{%
  \par\addvspace{0.5\baselineskip}%
  {\color{ealtwocolour}\ealtextblock{#1}}%
  \penalty200
  {\color{ealenglish}#2\par}%
  \addvspace{0.5\baselineskip}%
}

\begin{document}

% ---- Title ----
\begin{frame}[plain]
\begin{ealglossed}
  \centering
  {\LARGE\bfseries \ealgloss{Sequences}{Ciągi} Starters}\\[0.5em]
\end{ealglossed}
\end{frame}

% ================================================================
% STARTER 1
% ================================================================
\begin{frame}{Starter 1}
\begin{ealglossed}
  \begin{minipage}[t]{0.48\textwidth}
    \small\textbf{Q1.\enspace Matchstick Pattern}\\[0.5em]
    \matchsquares[red]{1}\quad
    \matchsquares[blue]{2}\quad
    \matchsquares[green!70!black]{3}\quad
    \ashowq{\matchsquares[orange]{4}}\\[0.4em]
    \footnotesize
    (a)\;Draw the 4th shape.\\
    (b)\;How many matchsticks does it use?\\[0.3em]
    \ashow{(b)\;13 matchsticks}
  \end{minipage}\hfill
  \begin{minipage}[t]{0.48\textwidth}
    \small\textbf{Q2.\enspace Complete the \ealgloss{Sequence}{ciąg}}\\[0.4em]
    \footnotesize
    5,\;10,\;15,\;20,\;
    \ablank{25},\;\ablank{30},\;\ablank{35}\\
    \ashow{(add 5 each time)}\\[0.6em]
    \small\textbf{Q3.\enspace Write the first 5 \ealgloss{terms}{wyrazów} of the \ealgloss{sequence}{ciągu} with \ealgloss{position-to-term rule}{wzór na n-ty wyraz ciągu}: $\boldsymbol{2n}$}\\[0.3em]
    \footnotesize
    \ablank{2},\;\ablank{4},\;\ablank{6},\;\ablank{8},\;\ablank{10}
  \end{minipage}
\end{ealglossed}
\end{frame}

\begin{frame}{Starter 1}
\begin{ealglossed}
  \small
  \textbf{Q4.\enspace}\;
  Sam writes the 3 \ealgloss{times table}{tabliczka mnożenia}: $3, 6, 9, 12, 15\ldots$ \\
  He then \textbf{adds 2} to every \ealgloss{term}{wyrazu}.
  Write the first 5 \ealgloss{terms}{wyrazów} of Sam's \ealgloss{sequence}{ciągu}.
  What is the \ealgloss{$n$th term}{wzór na n-ty wyraz} rule?\\[0.2em]
  \ashow{\ealgloss{Terms}{Wyrazy}:\;$5,\ 8,\ 11,\ 14,\ 17$.\quad
         \ealgloss{$n$th term}{wzór na n-ty wyraz}:\;$\boldsymbol{3n+2}$}
\end{ealglossed}
\end{frame}

\begin{frame}{Starter 1}
\begin{ealglossed}
  \begin{minipage}[t]{0.48\textwidth}
    \small\textbf{Q5.\enspace Extending Q1}\\
    \footnotesize
    Using the matchstick pattern from Q1:\\
    (a)\;How many matchsticks in the \textbf{10th} shape?\\
    (b)\;Write a rule for the $n$th shape.\\[0.2em]
    \ashow{(a)\;31 matchsticks.}\\
    \ashow{(b)\;\ealgloss{position-to-term rule}{wzór na n-ty wyraz ciągu}:\;$3n+1$}
  \end{minipage}\hfill
  \begin{minipage}[t]{0.48\textwidth}
    \small\textbf{Q6.\;Challenge}\\
    \footnotesize
    Is $\boldsymbol{100}$ a \ealgloss{term}{wyrazem} in the \ealgloss{sequence}{ciągu}
    with \ealgloss{$n$th term}{wzór na n-ty wyraz} $3n+1$\,?\\
    Show how you know.\\[0.2em]
    \ashow{Yes, when $n=33$,}\\
    \ashow{$3n+1=3\times33+1=99+1=100$}
  \end{minipage}
\end{ealglossed}
\end{frame}

% ================================================================
% STARTER 2
% ================================================================
\begin{frame}{Starter 2}
\begin{ealglossed}
  \begin{minipage}[t]{0.48\textwidth}
    \small\textbf{Q1.\enspace Dot Pattern}\\[0.5em]
    \fourtabledots[red]{1}\quad
    \fourtabledots[blue]{2}\quad
    \fourtabledots[green!70!black]{3}\quad
    \ashowq{\fourtabledots[orange]{4}}\\[0.4em]
    \footnotesize
    (a)\;Draw the 4th picture.\\
    (b)\;How many dots does it contain?\\[0.3em]
    \ashow{(b)\;16 dots}
  \end{minipage}\hfill
  \begin{minipage}[t]{0.48\textwidth}
    \small\textbf{Q2.\enspace Complete the \ealgloss{Sequence}{ciąg}}\\[0.4em]
    \footnotesize
    4,\;8,\;12,\;16,\;
    \ablank{20},\;\ablank{24},\;\ablank{28}\\
    \ashow{(add 4 each time)}\\[0.6em]
    \small\textbf{Q3.\enspace Write the first 5 \ealgloss{terms}{wyrazów} of $\boldsymbol{3n+1}$}\\[0.3em]
    \footnotesize
    \ablank{4},\;\ablank{7},\;\ablank{10},\;\ablank{13},\;\ablank{16}
  \end{minipage}
\end{ealglossed}
\end{frame}

\begin{frame}{Starter 2}
\begin{ealglossed}
  \small
  \textbf{Q4.\enspace}\;
  Alex writes the 5 \ealgloss{times table}{tabliczka mnożenia}: $5, 10, 15, 20, 25\ldots$
  He then \textbf{subtracts 1} from every \ealgloss{term}{wyrazu}.
  Write the first 5 \ealgloss{terms}{wyrazów} of Alex's \ealgloss{sequence}{ciągu}.
  What is the \ealgloss{$n$th term}{wzór na n-ty wyraz} rule?\\[0.2em]
  \ashow{\ealgloss{Terms}{Wyrazy}:\;$4,\ 9,\ 14,\ 19,\ 24$.\quad
         \ealgloss{$n$th term}{wzór na n-ty wyraz}:\;$\boldsymbol{5n-1}$}
\end{ealglossed}
\end{frame}

\begin{frame}{Starter 2}
\begin{ealglossed}
  \begin{minipage}[t]{0.48\textwidth}
    \small\textbf{Q5.\enspace Extend Q1}\\
    \footnotesize
    Using your dot pattern from Q1:\\
    (a)\;How many dots in the \textbf{10th} picture?\\
    (b)\;Write a rule for the $n$th picture.\\[0.2em]
    \ashow{(a)\;40 dots.\quad (b)\;Rule:\;$4n$}
  \end{minipage}\hfill
  \begin{minipage}[t]{0.48\textwidth}
    \small\textbf{Q6.\;Challenge}\\
    \footnotesize
    \ealgloss{Sequence}{Ciąg} A has \ealgloss{$n$th term}{wzór na n-ty wyraz} $4n-2$.\\
    \ealgloss{Sequence}{Ciąg} B has \ealgloss{$n$th term}{wzór na n-ty wyraz} $2n+6$.\\
    Do they share a \ealgloss{term}{wyraz} \\ \emph{at the same position}? \\
    \ashow{$4n-2=2n+6\;\Rightarrow\;n=4.$\;
           Both equal\;\textbf{14} at position $n=4$.}
  \end{minipage}
\end{ealglossed}
\end{frame}

% ================================================================
% STARTER 3
% ================================================================
\begin{frame}{Starter 3}
\begin{ealglossed}
  \begin{minipage}[t]{0.48\textwidth}
    \small\textbf{Q1.\enspace Matchstick Triangles}\\[0.5em]
    \matchtriangles[red]{1}\quad
    \matchtriangles[blue]{2}\quad
    \matchtriangles[green!70!black]{3}\quad
    \ashowq{\matchtriangles[orange]{4}}\\[0.4em]
    \footnotesize
    (a)\;Draw the 4th shape.\\
    (b)\;How many matchsticks does it use?\\[0.3em]
    \ashow{(b)\;12 matchsticks}
  \end{minipage}\hfill
  \begin{minipage}[t]{0.48\textwidth}
    \small\textbf{Q2.\enspace Complete the \ealgloss{Sequence}{ciąg}}\\[0.4em]
    \footnotesize
    7,\;14,\;21,\;28,\;
    \ablank{35},\;\ablank{42},\;\ablank{49}\\
    \ashow{(add 7 each time --- the 7$\times$ table)}\\[0.6em]
    \small\textbf{Q3.\enspace Write the first 5 \ealgloss{terms}{wyrazów} of $\boldsymbol{5n-2}$}\\[0.3em]
    \footnotesize
    \ablank{3},\;\ablank{8},\;\ablank{13},\;\ablank{18},\;\ablank{23}
  \end{minipage}
\end{ealglossed}
\end{frame}

\begin{frame}{Starter 3}
\begin{ealglossed}
  \small
  \textbf{Q4.\enspace}\;
  A \ealgloss{sequence}{Ciąg} \textbf{starts at 5} and \textbf{increases by 3} each time.
  Write the first 5 \ealgloss{terms}{wyrazów}.
  What is the \ealgloss{$n$th term}{wzór na n-ty wyraz} rule?\\[0.2em]
  \ashow{\ealgloss{Terms}{Wyrazy}:\;$5,\ 8,\ 11,\ 14,\ 17$.\quad
         \ealgloss{$n$th term}{wzór na n-ty wyraz}:\;$\boldsymbol{3n+2}$}
\end{ealglossed}
\end{frame}

\begin{frame}{Starter 3}
\begin{ealglossed}
  \begin{minipage}[t]{0.48\textwidth}
    \small\textbf{Q5.\enspace Extending Q1}\\
    \footnotesize
    Using your triangle pattern from Q1:\\
    (a)\;Write a rule for the $n$th shape.\\
    (b)\;How many matchsticks in the \textbf{20th} shape?\\[0.2em]
    \ashow{(a)\;Rule:\;$3n$.\quad (b)\;60 matchsticks.}
  \end{minipage}\hfill
  \begin{minipage}[t]{0.48\textwidth}
    \small\textbf{Q6.\;Challenge}\\
    \footnotesize
    The \ealgloss{$n$th term}{wzór na n-ty wyraz} of a \ealgloss{sequence}{ciągu} is $4n+3$.\\
    Which \ealgloss{term}{wyraz} is the \textbf{first to exceed 75}?\\
    Show your working clearly.\\[0.2em]
    \ashow{$4n+3>75\;\Rightarrow\;n>18.$\;
           19th \ealgloss{term}{wyraz}:\;$4(19)+3=\boldsymbol{79}.$}
  \end{minipage}
\end{ealglossed}
\end{frame}

% ================================================================
% STARTER 4
% ================================================================
\begin{frame}{Starter 4}
\begin{ealglossed}
  \begin{minipage}[t]{0.48\textwidth}
    \small\textbf{Q1.\enspace Dots}\\[0.5em]
    \rectdots[red]{1}\quad
    \rectdots[blue]{2}\quad
    \rectdots[green!70!black]{3}\quad
    \ashowq{\rectdots[orange]{4}}\\[0.4em]
    \footnotesize
    (a)\;Draw the 4th pattern.\\
    (b)\;How many dots does it contain?\\[0.3em]
    \ashow{(b)\;20 dots}
  \end{minipage}\hfill
  \begin{minipage}[t]{0.48\textwidth}
    \small\textbf{Q2.\enspace Complete the \ealgloss{Sequence}{ciąg}}\\[0.4em]
    \footnotesize
    1,\;5,\;9,\;13,\;
    \ablank{17},\;\ablank{21},\;\ablank{25}\\
    \ashow{(add 4 each time)}\\[0.6em]
    \small\textbf{Q3.\enspace Write the first 5 \ealgloss{terms}{wyrazów} of $\boldsymbol{10-2n}$}\\[0.3em]
    \footnotesize
    \ablank{8},\;\ablank{6},\;\ablank{4},\;\ablank{2},\;\ablank{0}
  \end{minipage}
\end{ealglossed}
\end{frame}

\begin{frame}{Starter 4}
\begin{ealglossed}
  \small
  \textbf{Q4.\enspace}\;
  Mia writes the 6 \ealgloss{times table}{tabliczka mnożenia}: $6, 12, 18, 24, 30\ldots$
  She then \textbf{subtracts 4} from every \ealgloss{term}{wyrazu}.
  Write the first 5 \ealgloss{terms}{wyrazów}, find the \ealgloss{$n$th term}{wzór na n-ty wyraz}, and find the 50th \ealgloss{term}{wyraz}.\\[0.2em]
  \ashow{\ealgloss{Terms}{Wyrazy}:\;$2,\ 8,\ 14,\ 20,\ 26$.\quad
         \ealgloss{$n$th term}{wzór na n-ty wyraz}:\;$6n-4$.\quad
         50th \ealgloss{term}{wyraz}:\;$\boldsymbol{296}$.}
\end{ealglossed}
\end{frame}

\begin{frame}{Starter 4}
\begin{ealglossed}
  \begin{minipage}[t]{0.48\textwidth}
    \small\textbf{Q5.\enspace Find the \ealgloss{$n$th Term}{wzór na n-ty wyraz}}\\
    \footnotesize
    $1,\;5,\;9,\;13,\;17,\;\ldots$\\[0.2em]
    What is the \ealgloss{$n$th term}{wzór na n-ty wyraz} rule?\\[0.2em]
    \ashow{\ealgloss{Differences}{Różnice}: $+4$. \;First \ealgloss{term}{wyraz}: 1.
           \quad Rule:\;$\boldsymbol{4n-3}$}
  \end{minipage}\hfill
  \begin{minipage}[t]{0.48\textwidth}
    \small\textbf{Q6.\;Challenge}\\
    \footnotesize
    The \ealgloss{$n$th term}{wzór na n-ty wyraz} of a \ealgloss{sequence}{ciągu} is $10-2n$.\\
    Which is the \textbf{last \ealgloss{positive}{liczba dodatnia} \ealgloss{term}{wyraz}}?\\
    What is the \textbf{first \ealgloss{negative}{ujemny} \ealgloss{term}{wyraz}}?\\
    \ashow{$n=4$:\;$+2$ (last \ealgloss{positive}{liczba dodatnia}).}\\
    \ashow{$n=5$:\;$0$.}\\
    \ashow{$n=6$:\;$-2$ (first \ealgloss{negative}{ujemny}).}
  \end{minipage}
\end{ealglossed}
\end{frame}

% ================================================================
% STARTER 5
% ================================================================
\begin{frame}{Starter 5}
\begin{ealglossed}
  \begin{minipage}[t]{0.48\textwidth}
    \small\textbf{Q1.\enspace Dots}\\[0.5em]
    \crossdots[red]{1}\quad
    \crossdots[blue]{2}\quad
    \crossdots[green!70!black]{3}\quad
    \ashowq{\crossdots[orange]{4}}\\[0.4em]
    \footnotesize
    (a)\;Draw the 4th shape.\\
    (b)\;How many dots in shape 5?\\[0.3em]
    \ashow{(b)\;Shape 5 has 17 dots}
  \end{minipage}
\end{ealglossed}
\end{frame}

\begin{frame}{Starter 5}
\begin{ealglossed}
  \small
  \textbf{Q5.\enspace}\;
  Jake says: \\ \textit{``I start at 3 and add 4 each time.''}
 \\ (a) Write his first 5 \ealgloss{terms}{wyrazów}, (b) find the \ealgloss{$n$th term}{wzór na n-ty wyraz} rule,
  and (c) decide whether $\boldsymbol{59}$ is in his \ealgloss{sequence}{ciągu}.
  Show your reasoning.\\[0.2em]
  \ashow{(a) $3,7,11,15,19$.}\\
  \ashow{(b) \ealgloss{$n$th term}{wzór na n-ty wyraz}:\;$4n-1$.}\\
  \ashow{(c) $4n-1=59\;\Rightarrow\;n=15.$\;\textbf{Yes!} (15th \ealgloss{term}{wyraz})}
\end{ealglossed}
\end{frame}

\begin{frame}{Starter 5}
\begin{ealglossed}
  \begin{minipage}[t]{0.48\textwidth}
    \small\textbf{Q2.\enspace Complete the \ealgloss{Sequence}{ciąg}}\\[0.4em]
    \footnotesize
    1,\;4,\;9,\;16,\;
    \ablank{25},\;\ablank{36},\;\ablank{49}\\
    \ashow{(\ealgloss{square numbers}{liczby kwadratowe},\;$n^2$)}\\[0.6em]
    \small\textbf{Q3.\enspace Write the first 5 \ealgloss{terms}{wyrazów} of $\boldsymbol{2n+5}$}\\[0.3em]
    \footnotesize
    \ablank{7},\;\ablank{9},\;\ablank{11},\;\ablank{13},\;\ablank{15}
  \end{minipage}
\end{ealglossed}
\end{frame}

\begin{frame}{Starter 5}
\begin{ealglossed}
  \begin{minipage}[t]{0.48\textwidth}
    \small\textbf{Q4.\enspace Find the \ealgloss{$n$th Term}{wzór na n-ty wyraz}}\\
    \footnotesize
    (a)\;$3,\;5,\;7,\;9,\;11,\;\ldots$\quad
    \ashow{$2n+1$}\\[0.3em]
    (b)\;$6,\;10,\;14,\;18,\;22,\;\ldots$\quad
    \ashow{$4n+2$}
  \end{minipage}
\end{ealglossed}
\end{frame}

\begin{frame}{Starter 5}
\begin{ealglossed}
  \begin{minipage}[t]{0.48\textwidth}
    \small\textbf{Q6.\;Challenge}\\
    \footnotesize
    Mia's \ealgloss{sequence}{ciąg} has \ealgloss{$n$th term}{wzór na n-ty wyraz} $3n+1$.\\
    Jake's \ealgloss{sequence}{ciąg} has \ealgloss{$n$th term}{wzór na n-ty wyraz} $2n+6$.\\
    Find the \textbf{first three values} that appear
    in \emph{both} \ealgloss{sequences}{ciągach}.
    What do you notice about them?\\[0.2em]
    \ashow{Shared values:\;$10,\;16,\;22,\;\ldots$
           \;(they increase by 6 each time).}
  \end{minipage}
\end{ealglossed}
\end{frame}

% ================================================================
% STARTER 6
% ================================================================
\begin{frame}{Starter 6}
\begin{ealglossed}
  \begin{minipage}[t]{0.48\textwidth}
    \small\textbf{Q1.\enspace Matchstick Triangles}\\[0.5em]
    \matchtriangles[red]{1}\quad
    \matchtriangles[blue]{2}\quad
    \matchtriangles[green!70!black]{3}\quad
    \ashowq{\matchtriangles[orange]{4}}\\[0.4em]
    \footnotesize
    (a)\;Draw the 4th shape.\\
    (b)\;How many matchsticks does it use?\\
    (c)\;Which \ealgloss{times tables}{tabliczki mnożenia} does the \ealgloss{sequence}{ciąg} represent?\\[0.3em]
    \ashow{(b)\;12 matchsticks}\\
    \ashow{(c)\;The 3 \ealgloss{times table}{tabliczka mnożenia}}
  \end{minipage}\hfill
  \begin{minipage}[t]{0.48\textwidth}
    \small\textbf{Q2.\enspace Complete the \ealgloss{Sequence}{ciąg}}\\[0.4em]
    \footnotesize
    4,\;8,\;12,\;16,\;
    \ablank{20},\;\ablank{24},\;\ablank{28}\\[0.5em]

    \small\textbf{Q3.\enspace The 3 \ealgloss{times table}{tabliczka mnożenia} has \ealgloss{nth term}{wzór na n-ty wyraz} rule $3n$, what is the \ealgloss{nth term}{wzór na n-ty wyraz} rule for the 5 \ealgloss{times table}{tabliczka mnożenia}?} \footnotesize
    \ablank{$5n$}\\[0.3em]
    
    \small\textbf{Q4.\enspace Write the first 5 \ealgloss{terms}{wyrazów} of $\boldsymbol{3n+4}$}\\[0.3em]
    \footnotesize
    \ablank{7},\;\ablank{10},\;\ablank{13},\;\ablank{16},\;\ablank{19}
  \end{minipage}
\end{ealglossed}
\end{frame}

\begin{frame}{Starter 6}
\begin{ealglossed}
  \small
  \textbf{Q5.\enspace}\;
  A \ealgloss{sequence}{Ciąg} is made by taking the \textbf{3 \ealgloss{times table}{tabliczka mnożenia}}
  and \textbf{subtracting 2 from every \ealgloss{term}{wyrazu}.}\\[0.1em]
  \footnotesize
  (a)\;Write the first 5 \ealgloss{terms}{wyrazów}.\\
  (b)\;What is the \ealgloss{$n$th term}{wzór na n-ty wyraz} rule?\\
  (c)\;Use your rule to find the \textbf{20th} \ealgloss{term}{wyraz}.\\
  (d)\;Is $\boldsymbol{100}$ in the \ealgloss{sequence}{ciągu}? Show your working clearly.\\[0.2em]
  \ashow{(a)\;1,\,4,\,7,\,10,\,13. \quad (b)\;$\boldsymbol{3n-2}$. \quad (c)\;$3\times20-2 = 58$.}\\
  \ashow{(d)\;$3n-2=100 \Rightarrow 3n=102 \Rightarrow n=34$. \textbf{Yes}, 100 is the 34th \ealgloss{term}{wyraz}.}
\end{ealglossed}
\end{frame}

% ================================================================
% STARTER 7
% ================================================================
\begin{frame}{Starter 7}
\begin{ealglossed}
  \begin{minipage}[t]{0.48\textwidth}
    \small\textbf{Q1.\enspace Dots}\\[0.5em]
    \fourtabledots[red]{1}\quad
    \fourtabledots[blue]{2}\quad
    \fourtabledots[green!70!black]{3}\quad
    \ashowq{\fourtabledots[orange]{4}}\\[0.4em]
    \footnotesize
    (a)\;Draw the 4th picture.\\
    (b)\;How many dots does it contain?\\
    (c)\;Which \ealgloss{times table}{tabliczka mnożenia} does the \ealgloss{sequence}{ciąg} represent?\\[0.3em]
    \ashow{(b)\;16 dots}\\
    \ashow{(c)\;The 4 \ealgloss{times table}{tabliczka mnożenia}}
  \end{minipage}\hfill
  \begin{minipage}[t]{0.48\textwidth}
    \small\textbf{Q2.\enspace Complete the \ealgloss{Sequence}{ciąg}}\\[0.4em]
    \footnotesize
    5,\;10,\;15,\;20,\;
    \ablank{25},\;\ablank{30},\;\ablank{35}\\
    \ashow{(add 5 each time --- the 5 \ealgloss{times table}{tabliczka mnożenia})}\\[0.6em]
    \small\textbf{Q3.\enspace Write the first 5 \ealgloss{terms}{wyrazów} of $\boldsymbol{4n-3}$}\\[0.3em]
    \footnotesize
    \ablank{1},\;\ablank{5},\;\ablank{9},\;\ablank{13},\;\ablank{17}
  \end{minipage}
\end{ealglossed}
\end{frame}

\begin{frame}{Starter 7}
\begin{ealglossed}
  \small
  \textbf{Q4.\enspace}\;
  Start with the \textbf{5 \ealgloss{times table}{tabliczka mnożenia}} and \textbf{add 1} to every \ealgloss{term}{wyrazu}.\\[0.1em]
  \footnotesize
  (a)\;Write the first 5 \ealgloss{terms}{wyrazów}.\\
  (b)\;What is the \ealgloss{$n$th term}{wzór na n-ty wyraz} rule?\\
  (c)\;Find the 20th \ealgloss{term}{wyraz}.\\
  (d)\;Is $\boldsymbol{104}$ in the \ealgloss{sequence}{ciągu}? Explain how you know.\\[0.2em]
  \ashow{(a)\;6,\,11,\,16,\,21,\,26. \quad (b)\;$\boldsymbol{5n+1}$. \quad (c)\;$5\times20+1 = 101$.}\\
  \ashow{(d)\;$5n+1=104 \Rightarrow 5n=103 \Rightarrow n=20.6$. \textbf{No}, 104 is not a \ealgloss{term}{wyraz}.}
\end{ealglossed}
\end{frame}

\end{document}